\subsection{MAXIMAL AND MINIMAL ELEMENTS}

DEFINITION 3. Let E be an ordered set. An element \(a \in E\) is said to be a minimal (resp. maximal) element of E if the relation \(x \leqslant a\) (resp. \(x \geqslant a\) ) implies \(x = a\) .

Every minimal element of E is a maximal element with respect to the opposite ordering, and vice versa.

Examples

\begin{enumerate}
\def\labelenumi{(\arabic{enumi})}
\item
  Let A be a set. In the subset of \(\mathfrak{P}(A)\) (ordered by inclusion) consisting of the non-empty subsets of A, the minimal elements are the subsets consisting of a single element.
\item
  In the set \(\Phi(\mathbf{E},\mathbf{F})\) of mappings of subsets of \(\mathbf{E}\) into \(\mathbf{F}\) (F being non-empty), ordered by the relation '' \(v\) extends \(u\) '' between \(u\) and \(v\) , the maximal elements are the mappings of the whole of \(\mathbf{E}\) into \(\mathbf{F}\) .
\end{enumerate}

* (3) In the set of natural integers \textgreater{} 1, ordered by the relation ``m divides n'' between m and n, the minimal elements are the prime numbers. *

* (4) The set of real numbers has no maximal element and no minimal element. *

